\begin{textsample}{NEG Dim 2 – 2012\_c5 – Score -2.00 – 2012\_c5/enade\_2012\_secr\_exec\_d3dc44a86b.txt}  \label{ex:f2_neg_078}
CAPÍTULO 5 DISTRIBUIÇÃO DOS CONCEITOS 5.1 PANORAMA NACIONAL DA DISTRIBUIÇÃO DOS CONCEITOS A Tabela 5.1 apresenta a quantidade e distribuição de cursos de Secretariado Executivo participantes do ENADE/2012, por faixa de conceito e Grande Região.
A diferença entre os cursos tabulados neste capítulo e no capítulo 2 são os cursos sem conceito, em princípio, aqueles sem alunos concluintes que participassem da prova ou que tivessem tão somente um aluno concluinte¹¹.
Observando -se os dados da Tabela 5.1, nota -se que, dos 74 cursos participantes, 30 ( 40,5 \% ) classificaram -se com conceito 3, o valor modal.
Este foi também o conceito modal nas regiões Norte ( 33,3 \% ), Nordeste ( 50,0 \% ) e Sul ( também 50,0 \% ).
A região Sudeste é bimodal, com os conceitos 3 e 4 ( 39,1 \%, correspondendo a nove cursos cada ).
Na região Centro-Oeste teve o conceito 2 como valor modal ( 54,5 \% ).
O conceito 4 foi o segundo mais frequente em nível nacional ( 21,6 \%, correspondendo a 16 cursos ) e o conceito 2, o terceiro ( 16,2 \%, correspondendo a 12 cursos ).
Houve, ainda, dez cursos ( 13,5 \% ) que receberam conceito 5 e outros dois cursos que receberam conceito 1 ( 2,7 \% ).
Apenas quatro dos 74 cursos de Secretariado Executivo ficaram sem conceito ( SC ).
A região Norte participou com seis cursos ou 8,1 \% do total nacional.
Nesta região, como já comentado, o conceito 3 foi atribuído a dois cursos dos seis participantes, o que equivale a 33,3 \% do total regional.
Dois cursos ficaram sem conceito ( SC ).
Os demais foram avaliados com conceito 1 ( um curso, correspondendo a 16,7 \% ), e conceito 5 ( também um curso ).
Nenhum curso foi alocado aos conceitos 2 ou 4. ¹¹ Estes cursos com somente um concluinte foram considerados no capítulo 2.
A região Nordeste participou com 14 cursos ou 18,9 \% do total nacional.
Como já comentado, destes, sete cursos, 50,0 \% em termos regionais, obtiveram conceito 3, o conceito modal para a região.
Os conceitos 2 e 4 foram atribuídos a, respectivamente, três e dois cursos ( 21,4 \% e 14,3 \% ).
Os conceitos 1 e 5 foram atribuído a um curso cada ( 7,1 \% ).
Nessa região tampouco, nenhum dos cursos ficou sem conceito ( SC ).
Como já comentado, dos 23 cursos participantes da região Sudeste, 18 obtiveram os conceitos 3 ou 4, os conceitos modais, com nove cursos ( 39,1 \% ) cada.
O conceito 2 foi atribuído a um curso ( 4,3 \% ) e o conceito 5, a três ( 13,0 \% ).
Apenas um curso ficou sem conceito ( SC ), enquanto nenhum foi alocado ao conceito 1.
A região Sul contou com 20 cursos ou 27,0 \% do total nacional.
Como já comentado, a predominância do conceito 3 foi de 50,0 \%, correspondentes a dez dos 20 cursos participantes na região Sul.
O conceito 2 foi atribuído a dois cursos ( 10,0 \% ), o conceito 4 a quatro cursos ( 20,0 \% ) e o conceito 5, a três deles ( 15,0 \% ).
Apenas um dos cursos da região Sul ficou sem conceito ( SC ), enquanto nenhum foi alocado ao conceito 1.
Como já comentado, seis dos 11 cursos participantes na região Centro-Oeste receberam conceito 2, o conceito modal ( 54,5 \% ).
Os demais cursos foram avaliados com conceito 3 ( dois cursos, 18,2 \% ), conceito 4 ( um curso, 9,1 \% ) e conceito 5 ( também dois cursos ).
Nesta região nenhum curso ficou sem conceito ( SC ) ou foi alocado no conceito 1. 5.2 CONCEITOS POR CATEGORIA ADMINISTRATIVA E POR GRANDE REGIÃO A Tabela 5.2 apresenta a distribuição dos cursos participantes do ENADE/2012 de Secretariado Executivo, por Categoria Administrativa, de acordo com os conceitos por eles alcançados, segundo as Grandes Regiões brasileiras.
Dos 74 cursos participantes, 21 ( 28,4 \% ) eram ministrados em instituições públicas e 53 ( 71,6 \% ), em privadas.
De acordo com as informações da Tabela 5.2, em termos nacionais, entre instituições públicas e privadas : das dez IES que receberam o conceito 5, sete eram públicas e três privadas.
Dos 21 cursos participantes de IES Públicas, os conceitos 3 e 5 foram os valores modais, atribuído a sete cursos cada ( 33,3 \% ).
Nesta categoria, somente um dos cursos ficou sem conceito ( SC ).
Entre os demais cursos de instituições públicas participantes, dois obtiveram conceito 2 ( 9,5 \% da categoria ) e quatro foram alocados ao conceito 4 ( 19,0 \% ).
Na rede privada, o conceito modal foi 3, com 23 cursos dos 53 da categoria.
Nesta categoria administrativa, três dos cursos ficaram sem conceito ( SC ).
Entre os demais cursos Na análise por região, observa -se que, na região Norte, as instituições privadas participaram com dois cursos ( 33,3 \% do total regional, a menor proporção entre todas as regiões ), dos quais um ficou sem conceito ( SC ) e o outro recebeu o conceito 1.
As instituições públicas participaram com quatro cursos ( 66,7 \% do total regional ), dos quais um ficou sem conceito ( SC ).
Dos demais cursos, dois foram alocados ao conceito 3 e o outro curso, ao conceito 5.
Nenhum curso foi alocado aos demais conceitos nesta categoria na região.
Na região Nordeste, a rede privada concentrou oito dos 14 cursos participantes, equivalentes a 57,1 \% do total da região.
Nesta Categoria Administrativa na região, \textbf{predominaram} os cursos com conceito 3, quatro cursos correspondendo a exata metade dos cursos oferecidos por IES privadas no Nordeste.
Os demais foram avaliados com conceitos 1 ( um curso correspondendo a 12,5 \% ), 2 ( dois cursos correspondendo a 25,0 \% ) e 4 ( também um curso ).
Nenhum curso recebeu o conceito 5.
As instituições públicas dessa região participaram com seis cursos ( 42,9 \% ), dos quais três obtiveram conceito 3, também o valor modal.
Os três demais foram avaliados com os conceitos 2, 4 e 5 ( um curso em cada ).
Nenhum curso ficou sem conceito ( SC ) nesta combinação de região e categoria.
Na região Sudeste, a proporção de cursos da rede privada, 95,7 \%, foi mais elevada do que nas demais regiões brasileiras, correspondendo a 22 dos 23 cursos participantes.
Nesta categoria, na região Sudeste, os conceitos modais foram 3 e 4 ( nove cursos cada ).
Nesta combinação de Categoria Administrativa e Grande Região, apenas um dos cursos ficou sem conceito ( SC ).
Os demais foram avaliados com conceito 2 ( um curso ) e conceito 5 ( dois cursos ).
As instituições públicas na região Sudeste participaram com somente um curso, o qual foi avaliado com conceito 5.
As instituições privadas concentraram 12 dos 20 cursos participantes da região Sul, 60,0 \% do total regional.
Nesta combinação de Categoria Administrativa e Grande Região, apenas um curso ficou sem receber conceito ( SC ).
Oito ficaram com conceito 3, o conceito modal.
Os demais foram avaliados com conceito 2 ( dois cursos ) e conceito 4 ( um curso ).
As instituições públicas na região Sul participaram com oito cursos ( 40,0 \% ), dos quais nenhum ficou sem conceito ( SC ).
O conceito modal ficou entre 4 e 5, com três cursos cada.
Além disso, o conceito 3 foi recebido por dois cursos.
Nenhum curso foi alocado aos conceitos 1 ou 2.
Na região Centro-Oeste, nove dos 11 cursos participantes eram de instituições privadas ( 81,8 \% em termos regionais ).
Destes, cinco concentraram -se no conceito 2, conceito modal.
Os demais receberam conceito 3 ( dois cursos ), conceito 4 ( um curso ) e conceito 5 ( também um curso ).
Nenhum curso ficou sem conceito ( SC ) nesta categoria, nem tampouco foi alocado ao conceito 1.
Dos dois cursos de instituições públicas, um foi avaliado no conceito 2 e o outro recebeu conceito 5.
Nesta região, nenhum curso ficou sem conceito ( SC ). 5.3 CONCEITOS POR ORGANIZAÇÃO ACADÊMICA E POR GRANDE REGIÃO Na Tabela 5.3 encontra -se a distribuição dos conceitos atribuídos aos cursos participantes do ENADE/2012 na Área de Secretariado Executivo, por Organização Acadêmica, segundo as Grandes Regiões brasileiras.
Dos 74 cursos participantes, 36 eram oferecidos em Universidades, 11 em Centros Universitários e os demais 27 em Faculdades.
Esta distribuição corresponde a, respectivamente, 48,6 \%, 14,9 \% e 36,5 \% dos cursos.
De acordo com os dados apresentados, de todos os dez cursos avaliados com conceito 5, nove eram vinculados a Universidades.
Apenas dois cursos em Universidades ficaram sem conceito ( SC ).
Esse tipo de Organização Acadêmica teve o conceito 3 como modal, com 14 cursos ( 38,9 \% ).
Os demais cursos de Universidades avaliados receberam os conceitos 2 ( três cursos ), 4 ( oito cursos ) e conceito 5 ( nove cursos, como já mencionado ).
Entre os cursos em Centros Universitários, o conceito modal também foi 3, com quatro cursos ( 36,4 \% ).
Neste tipo de organização acadêmica nenhum curso ficou sem conceito ( SC ) ou recebeu o conceito 5.
Os outros cursos neste tipo de Organização Acadêmica receberam os conceitos 1 ( dois cursos ), 2 ( outros dois cursos ) e conceito 4 ( três cursos ).
Nas Faculdades, apenas dois dos 27 cursos ficaram sem conceito ( SC ) e 12 ( 44,4 \% ) receberam, também, o mesmo conceito modal dos demais tipos de Organização Acadêmica, conceito 3.
Dos demais cursos neste tipo de Organização Acadêmica, sete receberam o conceito 2, cinco conceito 4 e um outro curso, o conceito 5.
Nenhum curso recebeu o conceito 1.
Considerando -se separadamente as regiões brasileiras, verifica -se que, na região Norte, as Universidades concentraram cinco dos seis cursos participantes.
Dos cursos em Universidades, dois ficaram sem conceito ( SC ), outros dois receberam conceito 3 e um curso recebeu conceito 5.
Os Centros Universitários da região Norte foram representados por somente um curso, o qual recebeu conceito 1.
As Faculdades não participaram com nenhum curso na região Norte.
Na região Nordeste, as Universidades participaram com seis dos 14 cursos na Área de Secretariado Executivo da região.
Nenhum dos cursos em Universidades no Nordeste ficou sem conceito ( SC ) ou recebeu conceitos 1 ou 2.
O conceito modal foi 3, com quatro cursos.
Os demais receberam conceitos 4 e 5 ( um curso cada ).
Os Centros Universitários contaram com somente um curso participante na região Nordeste, o qual foi avaliado com conceito 1.
As Faculdades foram representadas por sete cursos na região Nordeste.
Os conceitos modais foram 2 e 3, com três cursos cada, enquanto o curso restante recebeu conceito 4.
Na região Sudeste, as Universidades concentraram dez dos 23 cursos da região.
Entre os cursos em Universidades na região, o conceito modal foi 4 com quatro cursos, e nenhum curso ficou sem conceito ( SC ).
Os demais cursos receberam os conceitos 3 e 5, com três cursos cada.
Os Centros Universitários participaram com cinco cursos na região Sudeste, dos quais três obtiveram conceito modal, 3, e dois, o conceito 4.
As Faculdades foram representadas por oito cursos na região Sudeste, que se \textbf{distribuíram} nos conceitos 2 ( um curso ), 3 ( três cursos ) e 4 ( outros três cursos ), enquanto um curso ficou sem conceito ( SC ).
Dos 20 cursos da região Sul, 13 eram de Universidades, para os quais o conceito modal foi 3, com cinco cursos.
Nesse tipo de organização, nenhum dos cursos ficou sem conceito ( SC ) e os demais receberam os conceitos 2 ( dois cursos ), 4 ( três cursos ) e 5 ( também três cursos ).
Os Centros Universitários da região Sul contaram com dois cursos participantes, um com o conceito 3 e o outro com conceito 4.
Foram cinco os cursos vinculados a Faculdades na região Sul, um destes ficou sem conceito ( SC ) e os outros quatro ficaram com conceito 3.
Na região Centro-Oeste, dois dos 11 cursos eram de Universidades, um avaliado com conceito 2 e o outro, com conceito 5.
Os Centros Universitários da região Centro-Oeste contaram com dois cursos, ambos com conceito 3.
Dos sete cursos em Faculdades na região Centro-Oeste, nenhum ficou sem conceito ( SC ), três receberam conceito 2, dois conceito 3, um conceito 4 e o outro conceito 5.

% matched lemmas: distribuir, predominar
\end{textsample}
